% 第 4 章：移动储能（MESS）时空路由与 SOC
\section{移动储能时空路由与 SOC}\label{sec:res-mess}

\subsection{运输网络与行程}
MESS 在道路/电气图运输网络上路由（显式 \fld{TransportEdge}，:170，含 \fld{distance\_km}/
\fld{available}；或以电气图作代理 \fld{use\_electrical\_graph\_as\_transport\_proxy}）。从母线 $a$
到 $b$ 的行程时间
\begin{equation}
\tau^{travel}_{ab}=\frac{\mathrm{dist}_{ab}\ (\mathrm{km})}{\text{\fld{mess\_travel\_speed\_kmph}}\ (\mathrm{km/h})}.
\end{equation}
资源 $m$ 在 \fld{mobile\_storage\_available\_from\_hr}$[m]$ 前不得调度、成岛或出发；可用后可等待
并保留其后每个服务区间。

\subsection{SOC 与出力}
每个 MESS 单元的荷电状态按充放电效率演进
\begin{equation}
\mathrm{SOC}_{m,t+1}=\mathrm{SOC}_{m,t}+\big(\eta_{ch}P^{ch}_{m,t}-P^{dis}_{m,t}/\eta_{dis}\big)\Delta t,
\qquad 0\le\mathrm{SOC}_{m,t}\le \overline{E}_m,
\end{equation}
到达目标母线并连接后方可出力。逐步快照 \fld{MESSStateStep}（:274）记录 \fld{status}、
\fld{dispatch\_mw}、\fld{energy\_mwh}、\fld{soc}、\fld{arrival\_time\_hr}、\fld{remaining\_travel\_hr}、
\fld{target\_bus}。构网能力由单一字段 \srcpath{MobileStorage::grid\_forming} 决定；仅当该字段
与 \fld{allow\_mess\_black\_start} 同真时 MILP 允许 MESS 作根（黑启动）。

\subsection{三种装配口径}
结果 \fld{mess\_dispatch\_model}（:439）声明 MESS 出力如何得出：
\begin{itemize}
  \item \texttt{strict\_mip}：RA 阶段工作流对 RA 残余母线切负荷解专用 MESS 时空路由/SOC MILP
        （\fld{use\_strict\_mip\_for\_mess}=true）；
  \item \texttt{stage\_dispatch} / \texttt{stage\_dispatch\_fallback}：阶段内 MESS 调度，或在严格
        子问题不可解时回退（\fld{fallback\_to\_stage\_mess\_dispatch}）；
  \item \texttt{disabled}/\texttt{none}：未启用或无 MESS。
\end{itemize}
聚合量 \fld{mess\_energy\_delivered\_mwh} 与 \fld{mess\_travel\_distance\_km} 记于结果。
注意：在有限动作认证路径（第~\ref{sec:res-certified}~章）中，行程动态\textbf{不}在亚秒 DAE
时间轴上积分，到达与行程能量是主问题级前提。
